\chapter{The shaders}
\label{ch:shaders}

The GPU work is done by two short HLSL programs in \file{shaders/pbr.hlsl}. The program compiles
them at start-up and whenever the user presses F5, so students can edit the file while the
program runs. If a modified shader does not compile, the program keeps the previous version and
shows the compiler's messages in the Scene panel.

\section{The constant buffer}
All per-frame parameters travel in one constant buffer. HLSL packs constant buffers in rows of
16 bytes and does not let a vector straddle two rows, so the declaration is arranged in groups of
four scalars:

\excerpt{cbuffer}

The C++ structure that fills it must have exactly the same layout; a \code{static\_assert}
checks the most common mistake, a size that is not a multiple of 16:

\excerpt{frameconstants}

DirectXMath stores matrices by rows and multiplies row vectors on the left, $\vect p'=\vect pM$.
Declaring the HLSL matrices \code{row\_major} and writing \code{mul(p, M)} follows the same
convention, so no transposition is needed.

\section{Vertex shader: displacement}
\label{sec:vs}
The vertex shader is equation~\eqref{eq:displace}. It samples the displacement map at the mipmap
level \code{dispLod} chosen on the CPU (Chapter~\ref{ch:sampling}), moves the vertex along its
normal, and transforms position, normal and tangent to the world. It does \emph{not} recompute
the normal (Section~\ref{sec:displacement}).

\excerpt{vsmain}

\section{The bump normal on the GPU}
This is equation~\eqref{eq:nts} again, written to match the C++ version line by line:

\excerpt{bumpnormalts}

\section{Pixel shader: bump mapping and lighting}
The pixel shader proceeds in the order of the theory chapters:
\begin{enumerate}
\item Rebuild an orthonormal frame from the interpolated normal and tangent (Gram--Schmidt), and
the bitangent from the stored sign (Chapter~\ref{ch:tangent}).
\item Compute $\nts$ from the bump map and transform it with the TBN matrix.
\item Read the albedo (decoded from sRGB by the view) and the roughness, or the neutral values
chosen in the interface when a map is switched off.
\item Either output a debug view or evaluate the GGX/Cook--Torrance model of
Chapter~\ref{ch:lighting}, with the horizon factor.
\item Draw a magenta ring around the probed texel, in texture space.
\end{enumerate}

\excerpt{psmain}

Two details deserve a comment. The debug views write values such as $\frac12\vect n+\frac12$
that must reach the screen unchanged, but the render target encodes everything as sRGB; the
shader therefore applies the approximate inverse, $c^{2.2}$, beforehand. And the roughness is
clamped to at least $0.04$, because a perfectly smooth GGX surface has an infinitely thin
highlight that a point light cannot show.
